% !TEX root = ../main.tex
%% Corrigé du sujet type bac n° 6 — valeurs vérifiées : verifications/sujets_bac.py

\corrigeSujet{6}

\begin{corrigeBac}{9}{Le paracétamol}
\partieExo{1}{La synthèse}
\begin{questions}
\item Le paracétamol porte un groupe \res{\text{hydroxyle}}, lié au cycle (fonction phénol), et un groupe \res{\text{amide}}
$-\mathrm{NH}-\mathrm{CO}-$, caractéristique de la famille des \res{\text{amides}}. \pts{0,75 pt}
\item $n(\text{4-aminophénol})=\dfrac{3{,}0}{109{,}1}=\qty{2,75e-2}{mol}$ et $n(\text{anhydride})=\dfrac{4{,}0\times1{,}08}{102{,}1}=\qty{4,23e-2}{mol}$. Les réactifs
réagissent mole à mole : le \res{\text{4-aminophénol}} est limitant. \pts{1,25 pt}
\item $m_\mathrm{max}=\num{2.75e-2}\times151{,}2=\qty{4,16}{g}$ ; rendement $\eta=\dfrac{3{,}3}{4{,}16}=\res{79~\%}$. \pts{0,75 pt}
\item $\dfrac{151{,}2}{109{,}1+102{,}1}=\res{71{,}6~\%}$ : près de 30~\% de la masse des réactifs finit en acide éthanoïque. On peut valoriser ce
sous-produit, ou choisir une voie de synthèse qui n'en forme pas. \pts{1 pt}
\item La recristallisation \res{\text{purifie}} le solide : les impuretés restent dissoutes dans l'eau froide. La température de
fusion mesurée, \qty{169}{\degreeCelsius}, est celle du paracétamol, différente de celle du réactif (\qty{187}{\degreeCelsius}) : le produit est
du paracétamol pur. \pts{0,75 pt}
\item L'atome d'azote, plus électronégatif que C et H, porte un doublet non liant : c'est le \res{\text{site donneur}}. L'atome de
carbone du groupe C$=$O, lié à des atomes d'oxygène plus électronégatifs, porte une charge partielle positive : c'est le
\res{\text{site accepteur}}. \pts{0,75 pt}
\end{questions}
\partieExo{2}{Le contrôle d'un comprimé}
\begin{questions}[resume]
\item L'absorbance y est maximale : la mesure est plus \res{\text{sensible}} et moins affectée par un léger écart de réglage de la
longueur d'onde. \pts{0,5 pt}
\item $A=\varepsilon\,\ell\,c$ : l'absorbance est proportionnelle à la concentration. Les points sont alignés sur une droite passant par
l'origine : la loi est vérifiée, avec un coefficient directeur $k=\dfrac{0{,}777}{12{,}0}=\res{\qty{0,0648}{L.mg^{-1}}}$. \pts{1,25 pt}
\item $c=\dfrac{A}{k}=\dfrac{0{,}640}{0{,}0648}=\qty{9,88}{mg.L^{-1}}$ dans la solution diluée, donc \qty{494}{mg.L^{-1}} dans la fiole de \qty{1,000}{L} :
$\res{m=\qty{494}{mg}}$. $u(m)=0{,}015\times494=\qty{7,4}{mg}$ et $z=\dfrac{500-494}{7{,}4}=0{,}8<2$ : la masse est \res{\text{compatible}}
avec l'indication du fabricant. \pts{1,5 pt}
\item Sans dilution, l'absorbance vaudrait environ $0{,}0648\times494=32$ : bien au-delà du domaine d'étalonnage et des
possibilités du spectrophotomètre. \pts{0,5 pt}
\end{questions}
\schema[corrige]{bac-s6-etalonnage}
\end{corrigeBac}

\begin{corrigeBac}{6}{L'énergie d'un satellite}
\begin{questions}
\item Un métal éclairé émet des électrons seulement si la fréquence de la lumière dépasse un seuil, quelle que soit
l'intensité lumineuse. Le modèle ondulatoire ne l'explique pas ; dans le modèle du photon, chaque photon, d'énergie $h\nu$,
cède toute son énergie à un seul électron, qui n'est arraché que si $h\nu$ dépasse le travail d'extraction. \pts{1 pt}
\item $E=\dfrac{h\,c}{\lambda}=\dfrac{\num{6.63e-34}\times\num{3.00e8}}{\num{550e-9}}=\res{\qty{3,62e-19}{J}}$, soit $\res{\qty{2,26}{eV}}$. \pts{0,75 pt}
\item Un photon n'est absorbé que si $\dfrac{h\,c}{\lambda}\geqslant E_\mathrm{seuil}$, soit
\[\lambda_\mathrm{max}=\frac{h\,c}{E_\mathrm{seuil}}=\frac{\num{6.63e-34}\times\num{3.00e8}}{1{,}12\times\num{1.60e-19}}=\qty{1,11e-6}{m},\]
soit \res{\qty{1110}{nm}}, dans le \res{\text{proche infrarouge}}. \pts{1 pt}
\item $\dfrac{2{,}26-1{,}12}{2{,}26}=\res{50~\%}$ de l'énergie de ce photon est perdue sous forme thermique. \pts{0,75 pt}
\item $N=\dfrac{1361}{\num{3.62e-19}}=\res{\num{3.8e21}}$ photons par seconde et par mètre carré. \pts{0,75 pt}
\item En plein soleil, $P=0{,}30\times1361\times20=\qty{8,2}{kW}$, pendant $98{,}5-35=\qty{63,5}{min}$ : \res{\qty{8,6}{kW.h}} par orbite. Les
équipements consomment $5{,}0\times\dfrac{98{,}5}{60}=\res{\qty{8,2}{kW.h}}$ par orbite : la production suffit, avec une faible marge, à
condition de stocker $5{,}0\times\dfrac{35}{60}=\qty{2,9}{kW.h}$ dans les batteries pour l'éclipse, sans trop de pertes. \pts{1,75 pt}
\end{questions}
\end{corrigeBac}

\begin{corrigeBac}{5}{Le niveau sonore d'un concert}
\begin{questions}
\item $I=I_0\times10^{L/10}=\num{1.0e-12}\times10^{10{,}5}=\res{\qty{3,2e-2}{W.m^{-2}}}$. \pts{0,75 pt}
\item $P=I\times4\pi d^2=\num{3.16e-2}\times4\pi\times10^2=\res{\qty{40}{W}}$. \pts{0,75 pt}
\item $I$ est inversement proportionnelle à $d^2$ : quand $d$ double, $I$ est divisée par 4 et $L$ diminue de $10\log4=\qty{6}{dB}$. De
\qty{10}{m} à \qty{40}{m}, la distance double deux fois : $L=105-12=\res{\qty{93}{dB}}$. \pts{1 pt}
\item Les intensités s'ajoutent : l'intensité double et le niveau augmente de $10\log2=\qty{3}{dB}$ : $\res{\qty{96}{dB}}$. \pts{0,75 pt}
\item À \qty{96}{dB}, soit \qty{11}{dB} de plus que \qty{85}{dB}, la durée tolérée vaut $\dfrac{8\ \mathrm{h}}{2^{11/3}}\approx\res{\qty{38}{min}}$, bien moins que deux heures :
le spectateur s'expose à des \res{\text{lésions auditives}}. Avec les bouchons, le niveau tombe à \qty{71}{dB}, inférieur à \qty{85}{dB} :
le risque disparaît. \pts{1,25 pt}
\item Il faut perdre \qty{20}{dB}, soit diviser l'intensité par 100 et multiplier la distance par 10 : $\res{\qty{100}{m}}$. \pts{0,5 pt}
\end{questions}
\end{corrigeBac}
